\begin{frame}[t]{A sparse algorithm made 500-bp chromosome analysis feasible.}
\vspace{0pt}
{\fontsize{12.5}{15}\selectfont Chromosome 1 A/B compartment calculation at 500-bp resolution\par}
\begin{center}
\begin{tikzpicture}[x=1cm,y=.85cm,every node/.style={align=left,font=\fontsize{12.5}{15}\selectfont}]
\node[anchor=west,text=Teal] at (0,2.5) {\textbf{POSSUMM}\\Sparse matrix-vector operations};
\fill[Teal] (1.375,1.75) circle (.09);
\node[anchor=west,text=Teal,font=\fontsize{17}{20}\selectfont\bfseries] at (1.75,1.75) {23 GB RAM};
\node[anchor=west] at (7.2,2.05) {\textbf{2.5 min}\\Measured runtime};
\node[anchor=west,text=Muted] at (0,.75) {\textbf{Dense approach}\\Projected memory requirement};
\draw[Gold,line width=1.1pt] (10.12,.05)--(10.12,.3);
\draw[-{Latex[length=2mm]},Gold,line width=1.1pt] (10.12,.175)--(11.4,.175);
\node[anchor=east,text=Gold,font=\fontsize{17}{20}\selectfont\bfseries] at (9.85,.175) {$>4.6$ TB};
\draw[Ink,line width=.8pt] (0,-.55)--(11.4,-.55);
\foreach \xx/\lab in {0/10,3.8/100,7.6/{1{,}000},11.4/{10{,}000}} {
  \draw[Ink,line width=.8pt] (\xx,-.55)--(\xx,-.7);
  \node[anchor=north,align=center,font=\fontsize{10.5}{12.5}\selectfont] at (\xx,-.76) {\lab};
}
\node[anchor=north,align=center,text=Muted,font=\fontsize{10.5}{12.5}\selectfont] at (5.7,-1.25) {RAM in GB, logarithmic scale};
\end{tikzpicture}
\end{center}
\sources{Harris et al., including Eliaz, \href{https://www.nature.com/articles/s41467-023-38429-1}{Nature Communications 14:3303 (2023)}. Team result.}
\qadetail{Full primary source: Harris, Gu, Olshansky, Wang, Kaur, Eliaz et al., Chromatin alternates between A and B compartments at kilobase scale for subgenic organization, Nature Communications 14:3303 (2023), DOI 10.1038/s41467-023-38429-1. Yossi Eliaz is a coauthor, not the first author. Do not attribute implementation of POSSUMM solely to him without a contribution record. POSSUMM computes needed sparse matrix-vector products with a Lanczos-like method rather than materializing the dense correlation matrix. The slide's endpoint is the chromosome-1 A/B eigenvector at 500-bp resolution: measured 2.5 minutes and 23 GB RAM, versus the paper's projected greater-than-4.6-TB dense memory requirement. The greater-than marker expresses a reported lower bound, and the RAM axis is logarithmic. The paper does not report a measured dense runtime for this comparison. Separate endpoint: computing the first four principal components genome-wide at 500-bp resolution took 39 minutes and 77.65 GB RAM. The experiments sequenced more than 42 billion paired-end read pairs, which yielded 33 billion contacts after alignment, duplicate removal and quality control. The 33-billion figure is contacts, not raw paired-end reads. The experiments were completed in 2018 as an ENCODE phase-4 pilot; publication dates do not establish the dates of career transitions. Related coauthored genome-architecture work: Kaushal et al., CTCF loss has limited effects on global genome architecture in Drosophila despite critical regulatory functions, Nature Communications 12:1011 (2021). Related cancer-genomics work: Gilad, Eliaz et al., Communications Biology 4:399 (2021), DOI 10.1038/s42003-021-01929-1, used a genome-scale CRISPR dropout screen followed by drug-combination validation; six of eight tested combinations improved killing in MCF-7 cells, an in-vitro finding. The later reply and target-expression critique remain in the bibliography. Poolkeh, Eliaz, Danovich and Gasic, DOI 10.1101/2020.04.25.20079343, is a first-author 2020 medRxiv modeling preprint: SIR-D prevalence estimation plus nested-pool optimization produced a modeled Israeli example with pool sizes 92 then 10 and 295,951 tests for roughly nine million people, about 30-fold fewer than individual testing. It does not report a deployed population-screening or wet-lab experiment. These projects illustrate representation, computation and model assumptions, not empirical fork benefits.}
\note{[0:50] Moving into genomics made computational scale very concrete. Our chromosome study used thirty-three billion processed contacts. At five-hundred-base-pair resolution, a dense chromosome-one calculation was projected to need more than four point six terabytes of memory. POSSUMM used sparse matrix-vector operations instead. The published calculation took two and a half minutes and twenty-three gigabytes. The algorithm changed what we could compute before any choice of execution backend. For today's problem, that means identifying the useful computation and state first. The next challenge was running large numbers of analyses reproducibly.}
\end{frame}
